% Discussion of Figure 1 results (K=2 Pareto frontier).
% Intended to be \input{} from the main results section.

\paragraph{Setup.}
The $K{=}2$ portfolio comprises \texttt{Llama-3.1-8B} (budget,
\$0.00003/req) and \texttt{Gemini-2.5-Pro} (frontier, \$0.015/req),
spanning a $500{\times}$ cost ratio.
Quality is measured by DeepSeek-R1~\cite{deepseekai2025deepseekr1}
judgments on a scale of $[0,1]$;
\texttt{Gemini} scores 0.932 on average versus 0.793 for \texttt{Llama},
a 14-point gap.  Crucially, the weak model matches or exceeds the
strong model on approximately 40\% of prompts, providing substantial
headroom for contextual routing.
The bandit is trained on $n{=}8{,}374$ prompts and evaluated on a
held-out test set of $n{=}1{,}824$ prompts.  Hyperparameters
($\alpha{=}1.0$, $\neff{=}10$, hybrid
LinUCB~\cite{li2010contextual} with PCA-25 prompt
embeddings) are selected by Pareto AUC on a separate validation split
($n{=}1{,}785$).
All baseline hyperparameters are tuned on the same validation split
using the same criterion (Pareto AUC), ensuring a fair comparison
(Table~\ref{tab:tuned_hparams_k2}).

\paragraph{Reward signal assumption.}
\label{sec:reward_assumption}
\textit{banditGPT} requires per-request reward feedback but no
pre-collected routing labels.  In production, rewards can come from
implicit user signals---acceptance rates, thumbs-up/down ratings,
task-completion indicators---at no additional inference cost.
Our evaluation uses an LLM-as-judge
(DeepSeek-R1~\cite{deepseekai2025deepseekr1}) as a controlled
proxy for such signals, following the evaluation methodology
of~\citet{zheng2024llmjudge}.
Judge inference cost is an \emph{evaluation expense}, not a deployment
cost, and is therefore excluded from the cost axis in
Figure~\ref{fig:pareto_k2}.

\paragraph{Online bandits match supervised routing with weaker supervision.}
Figure~\ref{fig:pareto_k2} reveals a central finding:
\textit{banditGPT} achieves a Pareto AUC of
$0.870 \pm 0.001$ (5~seeds), matching the supervised logistic
regression baseline ($0.869$) despite receiving \emph{strictly less
information}.
The supervised router trains on full reward labels for \emph{all}
models on every prompt---the ``full information'' setting analogous
to RouteLLM~\cite{ong2024routellm} and RouterDC\@.
In contrast, the bandit observes only the reward for the model it
selected (partial, bandit feedback).
That the bandit closes this information gap demonstrates that
contextual exploration efficiently recovers the routing signal
that supervised methods extract from exhaustive labels.

\paragraph{Contextual features are necessary: context-free UCB collapses
to static mixing.}
Context-free UCB1 (Pareto AUC $0.862$) is statistically
indistinguishable from the static random-mix baseline ($0.863$).
This is expected: without prompt features, UCB1 converges to
a single global reward estimate per model.  Its selection rule
then reduces to a deterministic threshold on $\lambda$, producing
either ``always \texttt{Gemini}'' or ``always \texttt{Llama}''
with no intermediate operating points---identical to a biased
coin flip.
The ${\approx}0.8\%$ AUC gap between contextual methods ($0.870$)
and context-free UCB ($0.862$) confirms that prompt-level
features carry meaningful routing signal: certain prompts genuinely
require the frontier model while others are equally well-served by
the budget model.

\paragraph{Bandit algorithm choice is secondary for $K{=}2$.}
Among contextual methods---\textit{banditGPT} (LinUCB with warmup),
tabula rasa LinUCB, Linear Thompson Sampling, and
$\varepsilon$-greedy---Pareto AUC spans only $0.870$--$0.871$
($<0.15\%$ relative range).
This is not surprising: a two-arm portfolio presents a binary
routing decision that any reasonable contextual learner can master
given $n{=}8{,}374$ training examples.
The value of warmup priors, hybrid parameter sharing, and the
Corralling meta-learner---the architectural contributions of
\textit{banditGPT}---emerges in harder settings: $K{=}3$ with model
onboarding (Section~\ref{sec:onboarding}), limited training data,
and prior mismatch scenarios.

\paragraph{CostSave at matched quality.}
Table~\ref{tab:costsave_k2} quantifies the advantage at three
quality thresholds.  At each threshold $Q$,
CostSave~\cite{li2026llmrouterbench, chen2023frugalgpt} measures the
percentage reduction in cost required to achieve
$Q \times R_{\text{strong}}$ quality on the Pareto frontier, relative
to the cost of always using the strong model.

\begin{table}[htbp]
\centering
\small
\caption{CostSave\,@\,$Q$ for the $K{=}2$ portfolio
(\texttt{Llama-3.1-8B} vs.\ \texttt{Gemini-2.5-Pro}).
$\Delta$ is the bandit saving minus the static-baseline saving
at matched quality.  95\% CIs from prompt-level paired
bootstrap~\cite{efron1993bootstrap}
($n{=}1{,}824$, $B{=}2{,}000$).}
\label{tab:costsave_k2}
\begin{tabular}{lcccc}
\toprule
Threshold & Bandit & Static & $\Delta$ (pp) & 95\% CI \\
\midrule
CostSave\,@\,90\% & 71.9\% & 66.9\% & $+5.1$ & $[+2.0,\; +7.8]$ \\
CostSave\,@\,95\% & 41.5\% & 33.4\% & $+8.1$ & $[+5.1,\; +10.8]$ \\
CostSave\,@\,99\% & 14.2\% & \phantom{0}6.7\% & $+7.5$ & $[+4.5,\; +10.9]$ \\
\bottomrule
\end{tabular}
\end{table}

At the 95\% quality threshold---matching 95\% of \texttt{Gemini}'s
reward---the bandit saves 41.5\% of inference cost, an 8.1~pp
advantage over the static baseline ($95\%~\text{CI}:
[+5.1,\, +10.8]$~pp).  All three advantages exclude zero,
confirming that contextual routing provides a statistically
significant cost reduction beyond what non-contextual mixing achieves.

The advantage grows from 5.1~pp at $Q{=}90\%$ to 8.1~pp at
$Q{=}95\%$, reflecting the fact that moderate quality targets leave
the most room for the router to substitute the cheap model on
easy prompts.  At the stringent $Q{=}99\%$ threshold the absolute
saving shrinks (14.2\% vs.\ 71.9\% at $Q{=}90\%$) because nearly
all prompts must be routed to the frontier model, but the advantage
over the static baseline remains large ($+7.5$~pp).

\paragraph{Gap to oracle.}
\label{sec:gap_oracle}
The per-instance oracle---which always selects the higher-reward model
for each prompt---achieves an average reward of $0.946$.
We define the normalised Gap@Oracle as
$(R_{\text{oracle}} - R_{\text{router}}) /
 (R_{\text{oracle}} - R_{\text{weak}}) \times 100$,
where $0\%$ equals the oracle and $100\%$ equals always routing to the
weak model.
At the highest-quality operating point ($\lambda{=}0$), the bandit
reaches $R{=}0.926$ and a Gap@Oracle of \textbf{13.1\%}, closing
87\% of the weak-to-oracle quality range.
At moderate cost savings ($\lambda{=}0.20$), the gap widens to 35.3\%,
reflecting the deliberate trade-off of quality for cost.
Gap@Oracle is anchored to the portfolio's ceiling and floor rather than
to the cost ratio, making it more comparable across benchmarks than
CostSave alone.

\paragraph{Routing decisions are non-trivial.}
The routing mix at intermediate cost penalties reveals that the
bandit makes genuine discriminative decisions rather than trivially
defaulting to one model.
At $\lambda{=}0.20$, the router sends 42\% of prompts to
\texttt{Llama} and 58\% to \texttt{Gemini}, achieving a reward of
$0.892$ at \$0.0095/req---a 37\% cost reduction relative to
always-\texttt{Gemini} while retaining 96\% of its quality.
This confirms that the bandit learns a meaningful mapping from
prompt features to model selection, exploiting the 40\% of prompts
where \texttt{Llama} is adequate.
Such prompt-specific routing is the defining advantage of contextual
bandits over non-contextual
policies~\cite{li2010contextual, abbasi2011improved}.

\paragraph{Sensitivity to cost ratio and cross-benchmark comparison.}
The CostSave percentages in Table~\ref{tab:costsave_k2} are
specific to the $500{\times}$ cost ratio of this portfolio.
With a smaller cost gap, absolute savings would be lower;
conversely, a larger gap would amplify them.
The operationally meaningful quantity is the \emph{advantage}
$\Delta$ over the static baseline, which reflects the value of
contextual information independent of the cost scale.
Gap@Oracle (Section~\ref{sec:gap_oracle}) provides a complementary,
cost-ratio-invariant quality metric anchored to the per-instance
oracle ceiling.
Direct numerical comparison with CostSave figures from other
benchmarks (e.g., RouteLLM~\cite{ong2024routellm} on Chatbot
Arena~\cite{zheng2023judging}) is not valid, as the metric is
sensitive to the model pair, cost ratio, task distribution, and
evaluation methodology.
Like BARP~\cite{wang2025barp}, \textit{banditGPT} operates under
bandit feedback and produces a continuous cost--quality frontier.
Direct numerical comparison requires a shared benchmark; we leave
RouterBench~\cite{li2026llmrouterbench} evaluation to future work.

\paragraph{Statistical methodology.}
Confidence intervals use a prompt-level paired
bootstrap~\cite{efron1993bootstrap}
($n{=}1{,}824$ test prompts, $B{=}2{,}000$ resamples).
Per-prompt bandit outcomes are averaged across 5~seeds before
resampling.  The same resampled indices are applied to both the
bandit and static baselines, so the CI captures uncertainty in the
\emph{advantage} (difference), not the absolute saving.
The bandit continues learning during the test phase (online
updates), faithfully reflecting deployment
conditions~\cite{li2010contextual}.  After
$8{,}374$ training prompts the policy is near-convergent; test-phase
updates are incremental and the distinction from a frozen-policy
evaluation is negligible.

\paragraph{Practical implications.}
The Pareto frontier provides a deployment-ready
knob~\cite{wang2025barp, li2025llmbandit}: by adjusting
the cost penalty $\lambda$ at serving time, a practitioner can
continuously trade quality for cost.  For example:
\begin{itemize}[nosep,leftmargin=*]
    \item \emph{Quality-first} ($\lambda{=}0.05$): reward~$0.922$
          at \$0.0127/req (16\% cost reduction vs.\
          always-\texttt{Gemini}).
    \item \emph{Balanced} ($\lambda{=}0.25$): reward~$0.873$ at
          \$0.0078/req (49\% cost reduction).
    \item \emph{Budget-constrained} ($\lambda{=}0.50$): reward~$0.804$
          at \$0.0010/req (93\% cost reduction).
\end{itemize}
No retraining is required to move between these operating
points---the same trained bandit serves all of them.  This
flexibility is particularly valuable in production environments
where budget constraints fluctuate (e.g., peak vs.\ off-peak
traffic) or where different user tiers warrant different
quality--cost trade-offs~\cite{chen2023frugalgpt}.
